\section{Conclusioni e Sviluppi Futuri}
\label{sec:conclusions}

Il progetto \textbf{DashMe} dimostra come l'applicazione sistematica di tecniche di \textit{Business Intelligence}, modellazione ETL in Tableau e \textit{Visual Analytics} permetta di trasformare un archivio regolatorio complesso e testualmente eterogeneo come quello dei richiami FDA in un sistema decisionale leggibile e azionabile. Partendo dalle 17 colonne grezze di \texttt{data/recalls.xlsx} e costruendo 14 campi calcolati in \texttt{src/Tableau/dashsmash3.twb} su estrazione \texttt{data/dati\_trasformati\_tableau.hyper}, il lavoro consente di rispondere alle tre domande di ricerca formulate:

\begin{itemize}
    \item \textbf{Risposta a RQ1 (Geografia, settori e responsabilità aziendale)}: gli eventi di richiamo risultano fortemente concentrati in specifici poli produttivi e logistici statunitensi (fino a 4.522 eventi per singolo Stato, con preminenza di California, Texas, Florida, Nord-Est/Midwest e Puerto Rico). A livello settoriale, i comparti medicali, farmaceutici e cosmetici presentano una prevalenza di Classe~II (80,76\%--92,94\%, fino al 100,00\% nel tabacco), mentre \textit{Prodotti per Animali} (47,65\% Classe~I) e \textit{Cibo e Bevande} ($\approx$ 43,57\% Classe~I) concentrano la massima severità clinica. Nella graduatoria Top~15 per volume emergono grandi operatori medicali ed ematologici accanto ad aziende alimentari e farmaceutiche con elevatissima incidenza di Classe~I.
    \item \textbf{Risposta a RQ2 (Dinamiche temporali e tassonomia delle cause)}: l'architettura semantica per l'analisi temporale e causale è stata predisposta mediante i campi \texttt{Categoria Ragione di Richiamo}, \texttt{Dettaglio Allergene} e \texttt{Numero Allergeni Coinvolti}; la consolidazione empirica di trend, stagionalità e profili allergenici è formalmente strutturata nella Sezione~\ref{subsec:dash2} in attesa dell'esportazione definitiva di \texttt{Dash2.png}. \todo{Aggiornare la sintesi conclusiva di RQ2 con i risultati quantitativi estratti da Dashboard 2.}
    \item \textbf{Risposta a RQ3 (Governance operativa, catena del freddo, packaging e distribuzione)}: su 104.537 prodotti richiamati, 84.502 risultano terminati e 1.480 completati, ma ben 18.555 permangono ``In corso'', segnalando un'area critica di gestione post-notifica. Sul piano operativo, gli imballaggi in \textit{Plastica / Busta} ($>$8.000) e \textit{Vetro / Bottiglia / Barattolo} ($\approx$ 4.800) concentrano i richiami da \textit{Salmonella}, allergeni e \textit{Listeria}, mentre la matrice distributiva evidenzia come una quota significativa di eventi di Classe~I interessi reti di vendita regionali multi-stato e internazionali.
\end{itemize}

Il contributo principale di DashMe risiede nell'aver colmato il divario tra dato regolatorio grezzo e lettura orientata al \textit{Project Management} e al \textit{Risk Management}, offrendo una struttura analitica riproducibile per la prioritizzazione di audit, controlli di packaging e piani di ritiro.

In prospettiva futura, il progetto si presta a tre direttrici di evoluzione:
\begin{enumerate}
    \item \textbf{Finalizzazione della Dashboard~2}: completare l'esportazione grafica e la validazione quantitativa delle serie storiche, della stagionalità e della scomposizione multi-allergene;
    \item \textbf{Integrazione dinamica tramite API openFDA}: sostituire lo snapshot statico su foglio elettronico con una pipeline automatizzata verso gli endpoint ufficiali openFDA, abilitando l'aggiornamento incrementale dell'estrazione \texttt{.hyper} e il monitoraggio quasi in tempo reale dei richiami ``In corso'';
    \item \textbf{Modellazione predittiva della severità}: utilizzare le feature ingegnerizzate (settore merceologico, macro-causa, packaging, modalità di conservazione e portata distributiva) per addestrare modelli di classificazione capaci di stimare precocemente la probabilità di ricaduta in Classe~I al momento della segnalazione iniziale.
\end{enumerate}
